\section{Introduction}\label{sec:intro}

Uncontacted peoples --- in Brazil's official terminology, \emph{povos ind\'igenas isolados} --- are indigenous groups who maintain no sustained relations with the surrounding national society \citep{pringle2014,survival2025}. Brazil holds the world's largest number of such peoples: FUNAI, the National Foundation of Indigenous Peoples, maintained 114 official records of isolated peoples in 2024, of which the state had confirmed 29 and left twelve without any territorial protection instrument whatsoever \citep{funai_audit_2024}; a 2025 global survey estimates 124 groups in Brazil alone, roughly two-thirds of the world's 196 \citep{survival2025}. For these peoples, forest loss is not only an environmental problem: they depend entirely on the forest for food, medicine, and material culture, they have no acquired immunity to epidemic disease, and their isolation is itself a survival strategy --- formally protected since 1987 by FUNAI's no-contact policy \citep{pringle2014}. The epidemiological record underlines what is at stake: across 117 recorded contact-related epidemics in Greater Amazonia between 1875 and 2008, median mortality ranged between 18\% and 20\% per epidemic \citep{walker2015}.

A large and credible literature establishes that Indigenous Lands in the Brazilian Amazon, in general, inhibit deforestation and fire \citep{nepstad2006}, shelter over half of the region's remaining forest carbon \citep{soaresfilho2010}, and --- when backed by full collective property rights --- reduce clearing by around two-thirds relative to comparable boundaries \citep{baragwanath2020,blackman2017,baragwanath2023}. Roads and settlement expansion are consistently identified as the leading proximate drivers of pressure \citep{barber2014,andrade2021,ferrante2020}. Yet this literature treats Indigenous Lands as a single category. The lands that host \emph{uncontacted} peoples are different in three ways that matter for both science and policy: their inhabitants cannot be consulted, relocated, or compensated; the state's protection instrument for them (restrictions-of-use ordinances and exclusion belts, monitored by the Coordination for Isolated and Recently Contacted Peoples, CGiirc) is administratively fragile --- FUNAI's own 2024 internal audit found outdated registries, an authorised staffing deficit of 74 out of 156 field posts, and twelve records with no protection at all \citep{funai_audit_2024}; and any forest loss that reaches them carries existential, non-mitigable risk.

Despite this, no published study quantifies, at the level of individual territories and with matched controls, how much forest loss occurs inside and around the Indigenous Lands that host uncontacted peoples. Existing quantitative work either aggregates all Indigenous Lands \citep{nepstad2006,nolte2013}, works at basin scale across Indigenous Territories and Protected Areas jointly \citep{ritter2025}, or analyses specific road corridors \citep{ferrante2020}. Scholarship specifically on isolated peoples is overwhelmingly descriptive, legal, or advocacy-oriented; a 2025 global report documents threats qualitatively (logging, mining, agribusiness, missions) \citep{survival2025}, and Brazil's official record remains fragmented across FUNAI spreadsheets, civil-society registries, and protection ordinances \citep{funai_audit_2024,opi2025}. The gap is consequential: if uncontacted peoples' lands are, in effect, the strongest forest-buffers in the Amazon --- or, alternatively, quietly losing their buffering capacity along specific frontier segments --- nobody currently knows.

This paper closes that gap with a fully open, reproducible geospatial assessment. We assemble three public data layers --- official FUNAI territory boundaries, the mapped occupation areas and records of isolated peoples from the OPI observatory \citep{opi2025}, and the Hansen/UMD Global Forest Change dataset \citep{hansen2013,gfc_dataset} --- and measure forest-cover loss (2000--2023) inside 81 Indigenous Lands with isolated-peoples presence, against two complementary references: 155 area-matched non-isolated Indigenous Lands, and concentric buffer rings (0--10\,km and 10--50\,km) around each territory. We further relate surrounding pressure to accessibility (distance to mapped urban centers and infrastructure corridors) and quantify its concentration across territories. The study uses no field data, involves no contact of any kind, and publishes no settlement locations (Section~\ref{sec:ethics}).

The contributions are fourfold. \textbf{(C1)} A transparent, scripted data-integration pipeline that reconciles the official territory registry with the isolated-peoples record system and versioned forest-change rasters, released as open code. \textbf{(C2)} The first territory-level quantitative assessment of forest-loss pressure on Brazil's uncontacted-peoples lands, relative to matched controls and buffer rings, with effect sizes and uncertainty. \textbf{(C3)} A ranked, evidence-based exposure list identifying the specific territories where surrounding pressure is highest --- directly usable for protection prioritisation, and consistent with FUNAI's audited protection deficits. \textbf{(C4)} A documented, replicable ethics protocol for studying isolated peoples without endangering them.

\section{Related Work}\label{sec:related}

\subsection{Indigenous Lands and deforestation in the Amazon}
The protective effect of Indigenous Lands on Amazonian forests is among the best-documented findings in land-systems science. \citet{nepstad2006} showed that parks and Indigenous lands inhibited deforestation and fire across the Brazilian Amazon, with Indigenous lands --- covering one-fifth of the region --- acting as the region's most important barrier. \citet{soaresfilho2010} projected that the protected-area and Indigenous-land network, holding 54\% of remaining forests, would avoid large carbon emissions under business-as-usual scenarios. Methodological work, however, showed that naive comparisons overstate protection: protected lands are preferentially located in remote, low-pressure areas, and matched-control designs shrink the estimated effect \citep{andam2008,joppa2010}. \citet{nolte2013} found that strictly protected areas and Indigenous lands avoided the most deforestation under matching, while \citet{blackman2017} and \citet{baragwanath2020,baragwanath2023} established quasi-experimental evidence that formal tenure (titling in Peru; homologation in Brazil) substantially reduces clearing and can trigger regrowth. \citet{garnett2018} placed Indigenous lands globally at $\sim$38 million km$^2$ of managed or tenured land, and \citet{ritter2025} showed that across the Amazon basin, land within Indigenous Territories and Protected Areas is far less impacted by human activity (14--16\%) than unprotected land (38\%). None of these studies separates the lands hosting uncontacted peoples.

\subsection{Drivers of pressure: roads, small clearings, and frontier dynamics}
Roads are the canonical driver of Amazonian deforestation: \citet{barber2014} quantified the mitigating effect of protected areas along roads, and the BR-319 corridor literature documents how paving plans alone reorganise pressure on Indigenous lands, including those with isolated peoples \citep{ferrante2020,andrade2021}. \citet{kalamandeen2018} showed that small-scale clearing (patches below the 6.25\,ha reporting threshold of PRODES) now dominates Amazonian loss, motivating the use of the 30\,m Global Forest Change product rather than coarse official clearing inventories in this study.

\subsection{Uncontacted peoples: ethics, records, and protection}
Research involving uncontacted peoples is bounded by a strict non-contact norm codified in Brazilian policy since 1987 \citep{pringle2014} and in the international framework of ILO Convention 169 \citep{ilo169} and UNDRIP \citep{undrip}. Systematic knowledge therefore comes from indirect evidence: FUNAI's registry of records (114 entries in 2024, across three evidentiary categories) \citep{funai_audit_2024}, civil-society observatories such as OPI, whose platform mirrors the official categorisation (115 records: 29 confirmed, 26 under study, 60 information-only as of October 2026) \citep{opi2025}, and regional compilations \citep{raisg2024}. FUNAI's 2024 internal audit documents the administrative fragility of the protection system, including twelve records lacking any territorial protection instrument \citep{funai_audit_2024}. Quantitative, territory-level measurement of forest loss around these lands has been absent; this study supplies it without relaxing any ethical constraint.

